\documentclass[10pt,a4paper]{amsart}
\usepackage[margin=19mm]{geometry}
\usepackage{amsmath,amssymb,mathtools}
\usepackage[scr=rsfs,cal=euler]{mathalfa}
\usepackage{xfrac}
\usepackage[unicode,hidelinks,bookmarks=false]{hyperref}
\newcommand{\ie}{i.e.\ }
\newcommand{\cf}{cf.\ }
\newcommand{\resp}{respectively\ }
\newcommand{\Subsection}{\subsection}
\providecommand{\nobreakdash}{\nobreak\hbox{-}}
\setlength{\emergencystretch}{2em}
\multlinegap0pt
\usepackage{bm}
\usepackage{amssymb}
\usepackage{mathtools}
\usepackage[shortlabels]{enumitem}
\usepackage{algorithm}\usepackage{algpseudocode}
\newtheorem{lemma}{Lemma}
\newtheorem{theorem}{Theorem}
\newtheorem{corollary}{Corollary}
\title{Single-Chord Augmentation of Weighted Cycles for Algebraic Connectivity and Network Coherence}
\author{Jiarong Deng, Liu Chang, Quanshun Yang}
\date{Source: arXiv:2605.24479v1 (CC BY 4.0)}
\begin{document}
\maketitle

\noindent\emph{Source and licence.} Single-Chord Augmentation of Weighted Cycles for Algebraic Connectivity and Network Coherence. Version arXiv:2605.24479v1 (CC BY 4.0), distributed by arXiv under the Creative Commons Attribution 4.0 International licence, http://creativecommons.org/licenses/by/4.0/, as recorded in the arXiv licence metadata for this exact version and shown on the abstract page https://arxiv.org/abs/2605.24479v1. Full licence notice: \texttt{licenses/arxiv-2605.24479-CC-BY-4.0.txt}.\par\smallskip

\noindent\textit{Editorial scope.} Counted unit: the article's corollary cor:second-cosine (source0.tex lines 2002--2021). The article's complete original proof of it is delivered with it (source0.tex lines 2023--2182, one \texttt{proof} environment of 160 lines). No mathematics is written, repaired or re-derived: the statement and the proof are the article's own lines, in the article's own notation, and no retained line is substituted or re-flowed.  Retained uncounted as the source-local context the proof needs: lines 1716--1734; lines 1742--1764; lines 1755--1763; lines 1847--1854; lines 1883--1901; lines 1888--1896; lines 1898--1900.  Nothing was omitted from any retained span, so the package carries no omission record.  No retained line contains a citation, so the wrapper carries no bibliography and no bibliography entry.  No retained line contains a figure environment, an image-inclusion command or drawing code, so no image is delivered and the package declares no asset dependency.  Editorial formatting only, in addition: an AMS \texttt{amsart} preamble that restates the article's own declarations and macros used by the retained lines, namely the environments \texttt{lemma}, \texttt{theorem}, \texttt{corollary} on separate counters, together with the standard packages those lines need.  The retained environments print in this document as Lemma 1, Theorem 1, Corollary 1 in order of appearance, whereas the article prints the same items with the numbering of its own full document.  The complete native source of the article as distributed in the arXiv e-print for this exact version is bundled unchanged as \texttt{sources/201277/source0.tex}.
\begin{lemma}[Continuous spectral gap on the circle]\label{lem:circle_gap_disc}
	Let
	\[
	\mu_1:=\left(\frac{2\pi}{S}\right)^2,
	\qquad
	\mathcal U:=\operatorname{span}\left\{
	\sin\frac{2\pi t}{S},\ \cos\frac{2\pi t}{S}
	\right\}\subset H^1_{\mathrm{per}}(0,S).
	\]
	If $g\in H^1_{\mathrm{per}}(0,S)$ has mean zero and
	$g_\perp:=g-P_{\mathcal U}g$, then
	\begin{equation}\label{eq:circle_gap_disc}
		\int_0^S|g'(t)|^2\,\mathrm dt
		\ge
		\mu_1\int_0^S|g(t)|^2\,\mathrm dt
		+
		3\mu_1\int_0^S|g_\perp(t)|^2\,\mathrm dt.
	\end{equation}
\end{lemma}

\begin{lemma}[Low-energy mass transfer]\label{lem:low_energy_mass_transfer}
	Let $\boldsymbol{z}\in\mathbb{R}^n$ satisfy
	\[
	\boldsymbol{z}^\top L\boldsymbol{z}\le C_*n^{-2}\|\boldsymbol{z}\|_2^2.
	\]
	Then
	\begin{equation}\label{eq:low_energy_mass_transfer}
		\|f_{\boldsymbol{z}}\|_{L^2(0,S)}^2
		=
		\bar r\|\boldsymbol{z}\|_2^2\bigl(1+O_{\kappa,C_*}(\delta_n)\bigr).
	\end{equation}
	If $\boldsymbol{z}$ and $\boldsymbol{y}$ both satisfy the same
	low-energy bound, then
	\begin{equation}\label{eq:bilinear_mass_transfer}
		\left|
		\langle f_{\boldsymbol{z}},f_{\boldsymbol{y}}\rangle_{L^2(0,S)}
		-\bar r\langle\boldsymbol{z},\boldsymbol{y}\rangle_{\mathbb{R}^n}
		\right|
		\le
		C_{\kappa,C_*}\delta_n\bar r
		\|\boldsymbol{z}\|_2\|\boldsymbol{y}\|_2.
	\end{equation}
\end{lemma}

	\begin{equation}
		\left|
		\langle f_{\boldsymbol{z}},f_{\boldsymbol{y}}\rangle_{L^2(0,S)}
		-\bar r\langle\boldsymbol{z},\boldsymbol{y}\rangle_{\mathbb{R}^n}
		\right|
		\le
		C_{\kappa,C_*}\delta_n\bar r
		\|\boldsymbol{z}\|_2\|\boldsymbol{y}\|_2.
	\end{equation}

\begin{lemma}[Continuous mean control]\label{lem:continuous_mean_control}
	If $\boldsymbol{z}\perp\boldsymbol{1}$ and $f=f_{\boldsymbol{z}}$, then
	\begin{equation}\label{eq:continuous_mean_control}
		\left|\frac1S\int_0^Sf(t)\,\mathrm dt\right|
		\le
		C_\kappa\left(\Delta+\frac1n\right)S^{1/2}\|f'\|_{L^2(0,S)}.
	\end{equation}
\end{lemma}

\begin{theorem}[Fiedler vector is nearly sinusoidal in resistance arclength]\label{thm:fiedler-sine}
	There exists $\delta_0=\delta_0(\kappa)>0$ such that, if
	$\delta_n\le \delta_0$, then for every unit Fiedler vector
	$\boldsymbol{u}_1$ there exist a phase $\phi\in\mathbb{R}$ and a sign
	$\sigma\in\{\pm 1\}$ such that
	\begin{equation}\label{eq:fiedler_sine_sup}
		\max_{0\le i\le n-1}
		\left|
		u_{1,i}-\sigma\sqrt{\frac2n}
		\sin\left(\frac{2\pi s_i}{S}+\phi\right)
		\right|
		\le
		C_\kappa\sqrt{\frac{\delta_n}{n}}.
	\end{equation}
	Moreover,
	\begin{equation}\label{eq:lambda1_upper_for_later}
		\lambda_1\le \bar r\mu_1(1+C_\kappa\delta_n).
	\end{equation}
\end{theorem}

	\begin{equation}
		\max_{0\le i\le n-1}
		\left|
		u_{1,i}-\sigma\sqrt{\frac2n}
		\sin\left(\frac{2\pi s_i}{S}+\phi\right)
		\right|
		\le
		C_\kappa\sqrt{\frac{\delta_n}{n}}.
	\end{equation}

	\begin{equation}
		\lambda_1\le \bar r\mu_1(1+C_\kappa\delta_n).
	\end{equation}

\begin{corollary}[Second mode is nearly a cosine]\label{cor:second-cosine}
	Under the assumptions of Theorem~\ref{thm:fiedler-sine}, fix the unit
	Fiedler vector $\boldsymbol{u}_1$ and the phase $\phi$ from
	\eqref{eq:fiedler_sine_sup}. There exists a unit eigenvector
	$\boldsymbol{u}_2$ associated with $\lambda_2$, chosen orthogonal to
	$\boldsymbol{u}_1$, and a sign $\tau\in\{\pm1\}$ such that
	\begin{equation}\label{eq:second_cosine_sup}
		\max_{0\le i\le n-1}
		\left|
		u_{2,i}-\tau\sqrt{\frac2n}
		\cos\left(\frac{2\pi s_i}{S}+\phi\right)
		\right|
		\le
		C_\kappa\sqrt{\frac{\delta_n}{n}}.
	\end{equation}
	Moreover,
	\begin{equation}\label{eq:lambda2_upper_for_later}
		\lambda_2\le \bar r\mu_1(1+C_\kappa\delta_n).
	\end{equation}
\end{corollary}

\begin{proof}
	For each phase $\psi\in\mathbb{R}$, define
	\[
	x_i^{(\psi)}:=
	\sqrt{\frac2n}\sin\left(\frac{2\pi s_i}{S}+\psi\right),
	\qquad
	\widetilde{\boldsymbol{x}}^{(\psi)}:=P\boldsymbol{x}^{(\psi)}.
	\]
	The quadrature estimate used in the proof of
	Theorem~\ref{thm:fiedler-sine} is uniform in $\psi$, and yields
	\[
	\|\widetilde{\boldsymbol{x}}^{(\psi)}\|_2^2=1+O_\kappa(\delta_n),
	\qquad
	\frac{(\widetilde{\boldsymbol{x}}^{(\psi)})^\top
		L\widetilde{\boldsymbol{x}}^{(\psi)}}
	{\|\widetilde{\boldsymbol{x}}^{(\psi)}\|_2^2}
	\le
	\bar r\mu_1(1+C_\kappa\delta_n).
	\]
	Hence the two-dimensional space
	\[
	W:=
	\operatorname{span}\left\{
	P\left(\sqrt{\frac2n}\sin\frac{2\pi s_i}{S}\right)_{i=0}^{n-1},
	P\left(\sqrt{\frac2n}\cos\frac{2\pi s_i}{S}\right)_{i=0}^{n-1}
	\right\}
	\subset \boldsymbol{1}^\perp
	\]
	has maximal Rayleigh quotient at most
	$\bar r\mu_1(1+C_\kappa\delta_n)$. The min--max principle gives
	\eqref{eq:lambda2_upper_for_later}.
	
	Let $\boldsymbol{u}_2$ be a unit eigenvector associated with $\lambda_2$
	and orthogonal to $\boldsymbol{u}_1$. Set $f_j:=f_{\boldsymbol{u}_j}$
	for $j=1,2$. By Lemma~\ref{lem:low_energy_mass_transfer},
	\eqref{eq:lambda1_upper_for_later}, and \eqref{eq:lambda2_upper_for_later},
	\[
	\|f_j\|_{L^2(0,S)}^2=\bar r(1+O_\kappa(\delta_n)),
	\qquad j=1,2.
	\]
	For $j=1,2$, write
	\[
	g_j:=\frac{f_j}{\|f_j\|_{L^2(0,S)}},
	\qquad
	m_j:=\frac1S\int_0^S g_j(t)\,dt,
	\qquad
	\widehat g_j:=\frac{g_j-m_j}{\|g_j-m_j\|_{L^2(0,S)}}.
	\]
	Lemma~\ref{lem:continuous_mean_control} gives
	$|m_j|\le C_\kappa\delta_n S^{-1/2}$, hence
	$\|g_j-m_j\|_{L^2}=1+O_\kappa(\delta_n)$. Also,
	\[
	\int_0^S|\widehat g_2'(t)|^2\,dt\le \mu_1(1+C_\kappa\delta_n).
	\]
	By Lemma~\ref{lem:circle_gap_disc},
	\begin{equation}\label{eq:g2_close_U}
		\operatorname{dist}_{L^2(0,S)}(\widehat g_2,\mathcal U)
		\le C_\kappa\sqrt{\delta_n}.
	\end{equation}
	The proof of Theorem~\ref{thm:fiedler-sine} gives
	\begin{equation}\label{eq:g1_close_sine_for_cosine}
		\left\|
		\widehat g_1-\psi_{\sin}
		\right\|_{L^2(0,S)}
		\le
		C_\kappa\sqrt{\delta_n},
	\end{equation}
	where
	\[
	\psi_{\sin}(t):=
	\sigma\sqrt{\frac2S}\sin\left(\frac{2\pi t}{S}+\phi\right),
	\qquad
	\psi_{\cos}(t):=
	\sqrt{\frac2S}\cos\left(\frac{2\pi t}{S}+\phi\right).
	\]
	Then $\{\psi_{\sin},\psi_{\cos}\}$ is an orthonormal basis of
	$\mathcal U$.
	
	Since $\boldsymbol{u}_1\perp \boldsymbol{u}_2$, the bilinear transfer
	estimate \eqref{eq:bilinear_mass_transfer} implies
	\[
	\langle f_1,f_2\rangle_{L^2(0,S)}=O_\kappa(\delta_n\bar r),
	\]
	and therefore
	\[
	\langle g_1,g_2\rangle_{L^2(0,S)}=O_\kappa(\delta_n).
	\]
	Subtracting the small means and renormalizing changes the inner product by
	only $O_\kappa(\delta_n)$, so
	\[
	\langle \widehat g_1,\widehat g_2\rangle_{L^2(0,S)}
	=
	O_\kappa(\delta_n).
	\]
	Using \eqref{eq:g1_close_sine_for_cosine}, we obtain
	\begin{equation}\label{eq:g2_sine_projection_small}
		\left|
		\langle \widehat g_2,\psi_{\sin}\rangle_{L^2(0,S)}
		\right|
		\le
		C_\kappa\sqrt{\delta_n}.
	\end{equation}
	
	Decompose
	\[
	\widehat g_2=a\psi_{\sin}+b\psi_{\cos}+h_\perp,
	\qquad h_\perp\perp \mathcal U.
	\]
	By \eqref{eq:g2_sine_projection_small}, $|a|\le C_\kappa\sqrt{\delta_n}$.
	By \eqref{eq:g2_close_U}, $\|h_\perp\|_{L^2}\le C_\kappa\sqrt{\delta_n}$.
	Since $\|\widehat g_2\|_{L^2}=1$,
	\[
	b^2=1-a^2-\|h_\perp\|_{L^2}^2=1+O_\kappa(\delta_n).
	\]
	Choose $\tau\in\{\pm1\}$ so that $\tau b\ge 0$. Then
	\begin{equation}\label{eq:g2_L2_close_cosine}
		\|\widehat g_2-\tau\psi_{\cos}\|_{L^2(0,S)}
		\le C_\kappa\sqrt{\delta_n}.
	\end{equation}
	
	Also,
	\[
	\int_0^S|h_\perp'(t)|^2\,dt
	\le C_\kappa\mu_1\delta_n.
	\]
	Indeed,
	\[
	\int_0^S|h_\perp'|^2
	=
	\int_0^S|\widehat g_2'|^2-\mu_1(a^2+b^2)
	\le
	\mu_1(1+C_\kappa\delta_n)-\mu_1(1-\|h_\perp\|_2^2)
	\le
	C_\kappa\mu_1\delta_n.
	\]
	Now
	\[
	\widehat g_2-\tau\psi_{\cos}
	=
	a\psi_{\sin}+(b-\tau)\psi_{\cos}+h_\perp.
	\]
	Using $|a|\le C_\kappa\sqrt{\delta_n}$, $|b-\tau|\le C_\kappa\delta_n$,
	and the bound on $h_\perp'$, we obtain
	\[
	\|(\widehat g_2-\tau\psi_{\cos})'\|_{L^2(0,S)}
	\le
	C_\kappa\sqrt{\mu_1\delta_n}.
	\]
	Together with \eqref{eq:g2_L2_close_cosine}, the periodic Sobolev
	inequality yields
	\[
	\|\widehat g_2-\tau\psi_{\cos}\|_{L^\infty(0,S)}
	\le
	C_\kappa\sqrt{\frac{\delta_n}{S}}.
	\]
	Returning from $\widehat g_2$ to $g_2$ adds only
	$O_\kappa(\delta_n S^{-1/2})$. Multiplying by
	$\|f_2\|_{L^2(0,S)}=\sqrt{\bar r}(1+O_\kappa(\delta_n))$ and evaluating
	at $t=s_i$ proves \eqref{eq:second_cosine_sup}.
\end{proof}

\end{document}
